\documentclass[../../main.tex]{subfiles}
\begin{document}

{\bf [解答]}

$A = (a+bi)^p$ とおいて二項定理で展開すると，
\begin{align}
  A = (a^p - {}_pC_2 a^{p-2}b^2 + \cdots) + i({}_pC_1 a^{p-1}b - {}_pC_3 a^{p-3}b^3 + \cdots)
\end{align}
である．$a,b,p$は実数だから$A$の虚部は第二項であるから，$i$をのぞいたものを$B$とおくと，
\begin{align}
  B
  &= {}_pC_1 a^{p-1}b - {}_pC_3 a^{p-3}b^3 + {}_pC_5 a^{p-5}b^5 - \cdots \\
  &= \sum_{2k-1 \le p} (-1)^{k-1} {}_pC_{2k-1} a^{p-(2k-1)} b^{2k-1} \label{eq:1}
\end{align}
である．

題意を背理法で示す．$A \in \mathbb{R}$ つまり $B=0$ と仮定する．
素数$p$の値で場合分けして考える．

\subparagraph{$1^\circ$ $p=2$ の時}

この時は$A$を直接計算すると，
\begin{align}
  A = a^2 - b^2 + 2abi \quad \therefore B = 2ab
\end{align}
である．$a,b$が正の整数だから$B \neq 0$だがこれは仮定に矛盾する．

\subparagraph{$2^\circ$ $p \ge 3$ の時}

この時$p$ は素数ゆえ奇数である．\cref{eq:1}において
\begin{align}
  B & = {}_pC_1 a^{p-1}b - {}_pC_3 a^{p-3}b^3 + \cdots + (-1)^{(p-3)/2}{}_pC_{p-2}a^2b^{p-2} + (-1)^{(p-1)/2}{}_pC_p b^p \\
  & = (a\text{の倍数}) \pm b^p = 0
\end{align}
の形でかける．これは $b$ が $a$ の倍数であることを表し，$a \perp b$ に矛盾する．
\casesend

以上から題意は示された．

\end{document}
